% teorema-forcing-draft.tex -- Degrau 1: T2' as a forcing theorem.
% [DRAFT -- aguarda 15o parecer]
% Standalone; compiles with pdflatex.  Notation identical to 06-notas/paper/paper-v3.tex.
\documentclass[11pt]{article}
\usepackage[margin=1in]{geometry}
\usepackage{amsmath,amssymb,amsthm,mathtools}
\usepackage[T1]{fontenc}
\usepackage{microtype}
\usepackage[colorlinks=true,linkcolor=blue,citecolor=blue]{hyperref}
\usepackage{cleveref}

\newtheorem{theorem}{Theorem}
\newtheorem{lemma}{Lemma}
\newtheorem{proposition}{Proposition}
\theoremstyle{definition}
\newtheorem{definition}{Definition}
\theoremstyle{remark}
\newtheorem{remark}{Remark}
\crefname{lemma}{Lemma}{Lemmas}
\Crefname{lemma}{Lemma}{Lemmas}
\crefname{theorem}{Theorem}{Theorems}
\Crefname{theorem}{Theorem}{Theorems}

\newcommand{\R}{\mathbb{R}}
\newcommand{\hyb}{\mathrm{hyb}}
\newcommand{\thit}{t_{\mathrm{hit}}}
\newcommand{\texit}{t_{\mathrm{exit}}}
\newcommand{\trel}{t_{\mathrm{rel}}}
\newcommand{\tcross}{t_{\mathrm{cross}}}
\newcommand{\Ar}{\mathcal{A}}
\newcommand{\Tep}{T_{\mathrm{ep}}}
\newcommand{\res}{\mathrm{res}}
\DeclareMathOperator*{\esssup}{ess\,sup}

\title{The entry--exit law as a forcing theorem:\\
one dead unit driven by a general affine live subsystem}
\author{Leonardo Caloi Santos}
\date{2026-07-26}

\begin{document}
\maketitle

\begin{center}
\fbox{\parbox{0.9\linewidth}{\centering
\textbf{[DRAFT --- aguarda 15\textsuperscript{o} parecer]}\\[2pt]
\small This note is a full rewrite of Lemmas L1--L6 of \emph{paper-v3},
Appendix~A, with the residual $R$ treated as a \emph{prescribed forcing}
generated by an arbitrary live subsystem of the affine class, instead of by a
single live neuron. It has not been adjudicated. Two steps are declared open in
\cref{sec:open}; they are inherited from the published proof, not created here.}}
\end{center}

\section{Setting}\label{sec:setting}

\paragraph{Network and flow.} Data $(x_i,y_i)_{i=1}^n\in\R\times\R$ in general
position. Shallow network with \emph{fixed} readout weights $a_0,\dots,a_m$,
\begin{equation}\label{eq:net}
f_\theta(x)=a_0\,g_\beta(w_0x+b_0)+\sum_{k=1}^{m}a_k\,g_\beta(w_kx+b_k),
\qquad
g_\beta(z)=\beta^{-1}\log(1+e^{\beta z}),
\end{equation}
quadratic loss $L=\tfrac1{2n}\sum_i(f_\theta(x_i)-y_i)^2$ and gradient flow
$\dot\theta=-\nabla L$ on $[0,H]$. Unit $k=0$ is the unit $A$ that dies; we take
$a_0=+1$ (only $a_0$ enters $c$ and the margin equation, so this is a
normalization, not a restriction). The remaining $m$ units form the live
subsystem $B$, with parameters $\theta_B\in\R^{2m}$; $\theta_A=(w_0,b_0)$.
Set $\res_i=f_\theta(x_i)-y_i$.

\paragraph{Margin, residual, area.} Fix the datum $j$. The margin of $A$ on $j$
is $z=\theta_A\cdot(x_j,1)$, $c=\|(x_j,1)\|^2/n$, and the \emph{residual
excluding $A$} is
\begin{equation}\label{eq:residual}
R(t)=y_j-f_{-A}\bigl(x_j;\theta_B(t)\bigr),
\qquad
f_{-A}(x;\theta_B)=\sum_{k=1}^{m}a_k g_\beta(w_kx+b_k).
\end{equation}
The \emph{hybrid reference} $\theta^{\hyb}$ runs the ReLU flow until the first
time $\thit$ with $z=0$, then freezes $\theta_A$ and runs $\theta_B$ alone; write
$R^{\hyb}$ for its residual, $\Ar^{\hyb}(t)=\int_{\thit}^{t}R^{\hyb}(s)\,ds$,
$\Tep=\texit-\thit$, and $\theta^\infty$ for the hybrid limit of
Definition~1 of the paper (ReLU flow on $[0,\thit]$; $\theta_A$ frozen on
$[\thit,\texit]$; ReLU flow restarted at
$(\theta_A(\thit),\theta_B^{\hyb}(\texit))$ afterwards).

\paragraph{The live class.} We call $B$ \emph{affine of order $m$} on the window
if, on that window, the activation pattern of every live unit is constant and the
readout is fixed. Then $f_{-A}$ is \emph{linear in $\theta_B$}: with
$\varphi_i\in\R^{2m}$ the feature vector of datum $i$,
$(\varphi_i)_{2k-1}=a_k\mathbf 1_{k,i}x_i$, $(\varphi_i)_{2k}=a_k\mathbf
1_{k,i}$, one has $f_{-A}(x_i;\theta_B)=\varphi_i\cdot\theta_B$ up to the
softplus defect of \eqref{eq:softdefect} below, and the ReLU $B$-field is
$-G_B(\theta_B-\theta_B^\ast)$ with
\begin{equation}\label{eq:gram}
G_B=\frac1n\sum_{i=1}^n\varphi_i\varphi_i^{\!\top}\in\R^{2m\times2m},
\qquad G_B\succeq0 .
\end{equation}
This is the only structural fact about $B$ used below, and it holds for
\emph{every} $m$: a one-hidden-layer network with frozen activation pattern,
fixed readout and only the first layer trained is linear in its trainable
parameters. Note $G_B$ is a Gram matrix, hence positive semidefinite; the proof
of \cref{lem:gronwall} uses \emph{only} that sign, never $\lambda_{\max}(G_B)$.

We write $G$ for the full $(2m+2)\times(2m+2)$ Gram of the network with $A$ alive
and pattern frozen, and $G_{BA}$ for its $B\times A$ block; $G$, and hence its
principal submatrix $G_B$, is PSD.

\paragraph{Elementary bounds.} For all $z$,
\begin{equation}\label{eq:softdefect}
0\le g_\beta(z)-z_+\le\beta^{-1}e^{-\beta|z|},
\qquad
\bigl|\sigma(\beta z)-\mathbf 1_{\{z>0\}}\bigr|\le e^{-\beta|z|},
\qquad
\beta g_\beta(z)=\log(1+w),
\end{equation}
with $w=e^{\beta z}$, $u=1/w$, $\Psi(u)=\log u+u$, $\sigma(s)=1/(1+e^{-s})$, and
$\sigma(\beta z)=1/(1+u)$, $(1-\sigma(s))/\sigma(s)=e^{-s}$.

\section{Hypotheses}\label{sec:hyp}

All hypotheses are statements about the \emph{forcing} $R$ and about the
\emph{hybrid reference trajectory}, both computable by a single ReLU integration
with $A$ frozen. None of them constrains $m$, the dimension of $\theta_B$, or the
functional form of $f_{-A}$ beyond the affine class of \cref{sec:setting}.

\begin{itemize}
\item[(H0$'$)] \textbf{Regularity of the forcing.}
$R^{\hyb}\in C^1([\thit,H])$ with $\|\dot R^{\hyb}\|_\infty\le C_R$.

\item[(H1$'$)] \textbf{Transversal entry.} $R^{\hyb}(\thit)<0$; the arrival speed
is $v=c\,|R^{\hyb}(\thit)|>0$.

\item[(H2$'$)] \textbf{Transversal exit at the zero of the area.}
$\Ar^{\hyb}$ has a first zero at $\texit<H$, with $R^{\hyb}(\texit)>0$,
$\Ar^{\hyb}<0$ on $(\thit,\texit)$, and, for each $\delta>0$,
$|\Ar^{\hyb}|\ge a_0(\delta)>0$ on $[\thit+\delta,\texit-\delta]$.

\item[(H3$'$)] \textbf{Lipschitz response.} There is $\mathrm{Lip}_R$ with
$|R[\theta_B]-R[\theta_B']|\le\mathrm{Lip}_R\|\theta_B-\theta_B'\|$, and the
$B$-field is Lipschitz on a compact tube around the hybrid trajectory. In the
affine class this is automatic, with
$\mathrm{Lip}_R=\|\varphi_j\|$; we keep it as a hypothesis so that the statement
survives outside the class.

\item[(H4$'$)] \textbf{Isolated episode, uniform live margins (reinforced).}
No unit other than $A$ switches on $[0,H]$, and the live margins are
\emph{uniformly} bounded away from the kink:
\[
\bigl|z_k(x_i;\theta_B(t))\bigr|\ \ge\ \delta_B>0
\qquad\text{for all }k=1,\dots,m,\ i=1,\dots,n,\ t\in[0,H],
\]
along both the hybrid and the $\beta$-trajectory. The other margins of $A$ stay
$\le-\delta_1<0$.

\item[(H5$'$)] \textbf{Bounded residuals.} $|\res_i(t)|\le C_{\mathrm{res}}$ for
all $i$ and all $t\in[0,H]$, along both trajectories in the tube.
\end{itemize}

\begin{remark}[why (H4$'$) had to be reinforced]\label{rem:H4}
In the two-neuron case the single live unit contributes one softplus-versus-ReLU
discrepancy, absorbed silently into the $O(e^{-\beta\delta_1})$ bookkeeping of
(H3) of the paper. With $m$ live units there are $mn$ such discrepancies, and a
per-instance margin is not enough: the bound must be uniform over units, data and
time, or the sum of $mn$ terms is not controlled by a single exponential. The
uniform $\delta_B$ is what makes the new term of \cref{lem:gronwall}(step 2b)
beyond all orders in $\beta$. This is a genuine strengthening relative to (H4) of
the paper, not a restatement.
\end{remark}

\section{The theorem}\label{sec:thm}

\begin{theorem}[entry--exit law, forcing form]\label{thm:forcing}
Let $A$ be a softplus-$\beta$ unit with fixed readout $+1$ and margin $z$ on the
face $j$, coupled to a live subsystem $B$ affine of order $m$ in the sense of
\cref{sec:setting}, and let $R$ be the residual \eqref{eq:residual} it generates.
Assume \textup{(H0$'$)--(H5$'$)}. Then there is $\beta_0<\infty$ such that for
$\beta\ge\beta_0$:
\begin{enumerate}
\item[(i)] \emph{(sticking)} on compact subintervals of $(\thit,\texit)$ the
margin satisfies $u=c\beta|\Ar^{\hyb}(t)|(1+o(1))$, i.e.
\[
z_{\mathrm{park}}(t)=-\frac{\log\beta}{\beta}\bigl(1+o(1)\bigr);
\]
\item[(ii)] \emph{(release)} the first return of $z$ to $0$ after death obeys
\begin{equation}\label{eq:release}
t^\beta_{\mathrm{release}}
=\texit-\frac{1}{c\,R^{\hyb}(\texit)}\cdot\frac1\beta+o(\beta^{-1}),
\end{equation}
where $\texit$ is the first zero of the hybrid area $\Ar^{\hyb}$;
\item[(iii)] \emph{(uniform approximation)} for every $\delta>0$,
\begin{equation}\label{eq:sup}
\sup_{t\in[0,H]\setminus(B_\delta(\thit)\cup B_\delta(\texit))}
\bigl\|\theta^\beta(t)-\theta^\infty(t)\bigr\|=O_\delta\!\left(\frac{\log\beta}{\beta}\right),
\end{equation}
the $\log\beta$ being carried solely by the normal coordinate of $\theta_A$;
\item[(iv)] \emph{(live subsystem)}
$\sup_{[0,H]}\|\theta_B^\beta-\theta_B^{\hyb}\|\le M_2\log\beta/\beta^2$ with
\begin{equation}\label{eq:M2}
M_2=\frac{C_fC_4}{c},
\qquad
C_4=\frac{2}{c\min\bigl(|R^{\hyb}(\thit)|,R^{\hyb}(\texit)\bigr)}+\frac{\Tep}{a_0},
\qquad
C_f=\frac{\|\varphi_j\|}{n}.
\end{equation}
\end{enumerate}
\end{theorem}

Three things in this statement are new relative to Theorem~1 of the paper. First,
$B$ is any affine subsystem of any order $m$, not one live neuron. Second,
$M_2$ carries \emph{no} factor $e^{L_BH}$: the Gr\"onwall of
\cref{lem:gronwall} is replaced by an energy estimate that uses only
$G_B\succeq0$. Third, $C_f\le\|(x_j,1)\|\,(\sum_ka_k^2)^{1/2}/n$ grows at most
like $\sqrt m$, so every constant in \eqref{eq:M2} is polynomial in $m$; see
\cref{rem:uniform}.

\section{The six lemmas in forcing form}\label{sec:lemmas}

Each lemma below carries a \emph{provenance} line saying what is inherited from
Appendix~A of the paper and what is new. Where the provenance says
\emph{identical with $R$ generic}, we do not reproduce the proof: the audited
claim is that the published proof never uses the structure of $B$, only
(H0$'$)--(H2$'$), and reproducing it verbatim would add nothing. Where the
provenance says \emph{new}, the proof is written out in full.

\subsection*{L1$'$ --- normal form and additive quadrature}

\begin{lemma}[normal form and additive quadrature]\label{lem:gauge}
With $w=e^{\beta z}$, $u=1/w$, $\Psi(u)=\log u+u$ and $\tau=\beta t$,
\begin{equation}\label{eq:normalform}
\frac{dw}{d\tau}=c\,R(t)\,\frac{w^2}{1+w}\,(1+E_\beta),
\qquad |E_\beta|\le\frac{\log(1+w)}{\beta|R|},
\end{equation}
and one has the exact additive identity
\begin{equation}\label{eq:quadrature}
\frac{d\Psi(u)}{dt}=-c\beta R(t)+c\log(1+w)+\mathcal E(t),
\qquad
|\mathcal E(t)|\le C_3\beta e^{-\beta\delta_1}.
\end{equation}
Under \textup{(H1$'$)--(H2$'$)},
$\int_{\thit}^{\trel}c\log(1+w)\,ds\le C_4\log\beta/\beta$ with $C_4$ as in
\eqref{eq:M2}.
\end{lemma}

\noindent\emph{Provenance: identical to Lemma~1 of the paper with $R$ generic.}
Step~1 uses only $\res_j=g_\beta(z)-R$, which is the \emph{definition} of $R$,
and the scalar reduction (\cref{lem:corner-out}, step~4), which is a statement
about the other margins of $A$; steps~2--3 are algebra in $(\Psi,u,\sigma)$;
step~4 bounds the correction in three windows using only $R(\thit)\ne0$,
$R(\texit)\ne0$ and $|\Ar|\ge a_0$ in the bulk, and produces $C_4$ written purely
in terms of $R$. No line mentions $B$. The proof is unchanged; we record only
that $\Tep$, not $H$, is the correct length in the bulk term, since the bound is
taken over the episode.

\subsection*{L2$'$ --- the live subsystem: a PSD energy estimate}

This is the only lemma whose proof changes, and it changes in our favour.

\begin{lemma}[live subsystem and feedback, PSD form]\label{lem:gronwall}
Assume \textup{(H4$'$)}, \textup{(H5$'$)}, and, on $[\thit,s]$, the layer bound
$\int_{\thit}^{s}c\log(1+w)\,dt\le 2C_4\log\beta/\beta$ supplied by
\cref{lem:bootstrap}. Put $\Delta=\theta_B^\beta-\theta_B^{\hyb}$. Then
\begin{equation}\label{eq:Brate}
\sup_{[0,s]}\|\Delta\|
\;\le\;\underbrace{\frac{C_fC_4}{c}\cdot\frac{2\log\beta}{\beta^2}}_{\text{episode}}
\;+\;\underbrace{O(\beta^{-2})}_{\text{pre-hit}}
\;+\;\underbrace{H\Bigl(C_3e^{-\beta\delta_1}+C_5\sqrt m\,C_{\mathrm{res}}\,e^{-\beta\delta_B}\Bigr)}_{\text{beyond all orders}} ,
\end{equation}
so $\sup\|\Delta\|=O(\log\beta/\beta^2)$ with constant $M_2=C_fC_4/c$, and
consequently
\begin{equation}\label{eq:Rpert}
R^\beta=R^{\hyb}+\delta,\qquad
\|\delta\|_\infty\le\mathrm{Lip}_R\sup\|\Delta\|
+\frac{m}{\beta}e^{-\beta\delta_B} .
\end{equation}
\emph{No factor $e^{L_BH}$ occurs, and $\lambda_{\max}(G_B)$ does not appear.}
\end{lemma}

\noindent\emph{Provenance: new proof.} The published Lemma~2 obtains the same
rate through the integral Gr\"onwall
$\|\Delta(t)\|\le\int e^{L_B(t-s)}\|G_A(s)\|ds$, paying $e^{L_BH}$ with
$L_B=\lambda_{\max}(G_B)$; with unit readouts $\lambda_{\max}(G_B)$ grows like
$m$, so that route makes $M_2$ exponential in $m$ and the theorem vacuous for
large $m$. The observation below is that the sign of $G_B$ is available and
Gr\"onwall throws it away. Step~2b (the live-unit discrepancies) has no
counterpart in the published proof.

\begin{proof}
\emph{Step 1 (exact perturbed linear system on the episode).} On $[\thit,H]$ the
hybrid freezes $\theta_A$ at the value with $z=0$, so $A$'s contribution to the
hybrid residuals is $\mathrm{ReLU}(z)\equiv0$ on the face $j$ and
$O(e^{-\beta\delta_1})$ elsewhere, and the hybrid $B$-flow is exactly
$\dot\theta_B^{\hyb}=-G_B\theta_B^{\hyb}+b$. In the $\beta$-flow the same
computation applies with two extra sources: the output of $A$, which enters every
residual, and the softplus defects of the live units themselves. Hence
\begin{equation}\label{eq:Deltaode}
\dot\Delta=-G_B\Delta+G_A(t)+G_S(t),
\end{equation}
where, with $\varphi_i$ the feature vectors of \eqref{eq:gram},
\[
G_A(t)=-\frac1n\sum_{i=1}^n g_\beta\bigl(z_{A,i}(t)\bigr)\varphi_i,
\qquad
G_S(t)=\text{(live softplus defect, bounded in step 2b).}
\]
Equation \eqref{eq:Deltaode} is an identity, not an estimate: the affine class
makes the drift exactly linear with the \emph{same} matrix $G_B$ for both flows,
which is why the difference of the drifts is $-G_B\Delta$ with no remainder.

\emph{Step 2a (bound on $G_A$).} $\|\varphi_j\|/n=:C_f$. By (H4$'$) the margins
$z_{A,i}$, $i\ne j$, are $\le-\delta_1$, so those terms contribute at most
$C_3e^{-\beta\delta_1}$; the term $i=j$ is $C_fg_\beta(z)$. Hence
\begin{equation}\label{eq:GA}
\|G_A(t)\|\le C_f\,g_\beta(z(t))+C_3e^{-\beta\delta_1}
=\frac{C_f}{\beta}\log\bigl(1+w(t)\bigr)+C_3e^{-\beta\delta_1}.
\end{equation}
Nothing in \eqref{eq:GA} depends on $m$ except through $C_f$: the size of the
perturbation that $A$ injects is a property of $A$.

\emph{Step 2b (the live-unit discrepancies; new term).} The $\beta$-flow's
$B$-field is
$-\tfrac1n\sum_i\res_i\,a_k\sigma(\beta z_{k,i})(x_i,1)$ per unit $k$, whereas
the hybrid uses $a_k\mathbf 1_{\{z_{k,i}>0\}}(x_i,1)$; and $\res_i$ itself
contains $\sum_ka_kg_\beta(z_{k,i})$ where the hybrid has
$\sum_ka_k(z_{k,i})_+$. By (H4$'$) every live margin satisfies
$|z_{k,i}|\ge\delta_B$, so \eqref{eq:softdefect} gives, per unit and per datum,
a derivative defect $\le e^{-\beta\delta_B}$ and an output defect
$\le\beta^{-1}e^{-\beta\delta_B}$. Using (H5$'$) to bound $|\res_i|\le
C_{\mathrm{res}}$ and $\|(x_i,1)\|\le X_{\max}$, the $2m$-vector $G_S$ obeys
\begin{equation}\label{eq:GS}
\|G_S(t)\|\;\le\;
\Bigl(\sum_{k=1}^m a_k^2\Bigr)^{1/2}X_{\max}
\Bigl(C_{\mathrm{res}}+\tfrac{m}{\beta}\max_k|a_k|X_{\max}\Bigr)e^{-\beta\delta_B}
\;=:\;C_5\sqrt m\,C_{\mathrm{res}}\,e^{-\beta\delta_B}\,(1+o(1)),
\end{equation}
the crude count being $mn$ scalar discrepancies, each $O(e^{-\beta\delta_B})$,
recombined into a Euclidean norm. The dependence on $m$ is at worst $m$ (at worst
$\sqrt m$ from the readout norm times $\sqrt m$ from the second bracket) and
multiplies a quantity beyond all orders in $\beta$; for
$\beta\ge\beta_1:=\delta_B^{-1}\log(m^2 H\beta^2)$ the whole term is
$o(\beta^{-2})$.

\emph{Step 3 (energy estimate; this is where PSD replaces Gr\"onwall).}
Let $\rho(t)=\|\Delta(t)\|$. From \eqref{eq:Deltaode},
\[
\frac{d}{dt}\tfrac12\rho^2
=-\Delta^{\!\top}G_B\Delta+\Delta^{\!\top}(G_A+G_S)
\;\le\;\rho\,\bigl(\|G_A\|+\|G_S\|\bigr),
\]
because $G_B\succeq0$ by \eqref{eq:gram}. Wherever $\rho>0$ this gives
$\dot\rho\le\|G_A\|+\|G_S\|$; since $\rho$ is Lipschitz, the same inequality
holds for its upper Dini derivative everywhere, and comparison with
$y'=\|G_A\|+\|G_S\|$, $y(\thit)=\rho(\thit)$, yields
\begin{equation}\label{eq:psd}
\rho(t)\ \le\ \rho(\thit)+\int_{\thit}^{t}\bigl(\|G_A\|+\|G_S\|\bigr)ds .
\end{equation}
Substituting \eqref{eq:GA} and using the layer bound of the hypothesis,
\[
\int_{\thit}^{s}\|G_A\|\,ds
\le\frac{C_f}{c\beta}\int_{\thit}^{s}c\log(1+w)\,ds+HC_3e^{-\beta\delta_1}
\le\frac{C_fC_4}{c}\cdot\frac{2\log\beta}{\beta^{2}}+HC_3e^{-\beta\delta_1},
\]
which is the first and third groups of \eqref{eq:Brate}, with
$M_2=C_fC_4/c$.

\emph{Step 4 (the pre-hit initial gap $\rho(\thit)$).} On $[0,\thit]$ both flows
have $A$ alive with frozen pattern, so the \emph{full} system is affine with the
full Gram $G\succeq0$ and $\delta\theta=\theta^\beta-\theta^{\mathrm{ReLU}}$
obeys $\dot{\delta\theta}=-G\,\delta\theta+E$, where $E$ collects the softplus
defects. Repeating step~3 on the full vector,
$\|\delta\theta(t)\|\le\int_0^t\|E\|$. The $A$-block of $E$ is
$O(e^{-\beta z^0(t)})$ (a derivative defect, no $\beta^{-1}$), and near the hit
transversality gives $z^0(t)=v(\thit-t)(1+O(\delta_z))$, so
\begin{equation}\label{eq:prehitA}
\|\delta\theta_A(t)\|\ \le\ \frac{C}{\beta v}\,e^{-\beta z^0(t)}\bigl(1+o(1)\bigr),
\end{equation}
which is step~1 of \cref{lem:corner-in} in pointwise form; it integrates to
$O(\beta^{-1})$ at $t=\thit$. The $B$-block of $E$ contains \emph{no} derivative
defect of $A$ --- $B$'s gradients involve only the live units' $\sigma$'s --- but
only $A$'s output defect $\le C_f\beta^{-1}e^{-\beta z^0}$ and $G_S$. Splitting
the $B$-block of the drift as $-G_B\Delta-G_{BA}\delta\theta_A$ and repeating
step~3 on $\Delta$ alone (again using $G_B\succeq0$; the cross term is treated as
a source),
\[
\rho(\thit)\le\int_0^{\thit}\Bigl(\|G_{BA}\|\,\|\delta\theta_A\|
+\frac{C_f}{\beta}e^{-\beta z^0}+\|G_S\|\Bigr)ds
\le\frac{C\|G_{BA}\|}{(\beta v)^2}+\frac{C_f}{\beta^2v}+HC_5\sqrt mC_{\mathrm{res}}e^{-\beta\delta_B},
\]
using $\int_0^{\thit}e^{-\beta v(\thit-s)}ds=(\beta v)^{-1}(1+o(1))$ twice. Hence
$\rho(\thit)=O(\beta^{-2})$, which is the second group of \eqref{eq:Brate} and is
dominated by the episode term. Estimate \eqref{eq:Rpert} follows from (H3$'$)
plus the $m\beta^{-1}e^{-\beta\delta_B}$ discrepancy between $R^\beta$ and
$R^{\hyb}$ at equal $\theta_B$.
\end{proof}

\begin{remark}[what the PSD step buys]\label{rem:psd}
The published route pays $e^{L_BH}$ because Gr\"onwall's inequality is stated for
$\|\dot\Delta\|\le L_B\|\Delta\|+\|G_A\|$ and is blind to the fact that the
linear part is \emph{dissipative}. Here $-G_B$ contracts, so the homogeneous part
can only help; \eqref{eq:psd} is an equality-free statement of that. The price is
that the argument is confined to the affine class: outside it, $-G_B\Delta$
acquires a remainder and one is back to a Lipschitz constant. Within the class,
which is exactly the class the theorem is stated for, the improvement is
unconditional and removes the only exponential in $m$ from the constants.
\end{remark}

\subsection*{L3$'$ --- entry corner}

\begin{lemma}[entry corner]\label{lem:corner-in}
Let $v=c|R^{\hyb}(\thit)|>0$ (transversal by (H1$'$)). Then
\begin{equation}\label{eq:corner-in}
\tcross^\beta=\thit+\frac{1}{\beta v}\bigl(1+O(\delta_z)+o(1)\bigr),
\qquad
\|\theta^\beta(\tcross)-\theta^\infty(\thit)\|=O(\beta^{-1}),
\end{equation}
$\tcross^\beta$ being the first time with $w=w_0$; the exponentially weighted
approach integral equals $\tfrac{3/2}{\beta v}(1+o(1))$, the $3/2$ coming from
$\int_0^\infty e^{-s}(1+e^{-s})ds$.
\end{lemma}

\noindent\emph{Provenance: identical to Lemma~3 of the paper with $R$ generic.}
Both steps are scalar in $z$: step~1 is the monotone sub-solution iteration
against the weight $e^{-\beta z^0}$, step~2 is the window-crossing identity
$(1-\sigma)/\sigma=e^{-s}$. $R$ enters only through $v$, i.e. only through
(H1$'$); $B$ appears zero times. We used the pointwise form of step~1 in
\eqref{eq:prehitA}, which is the same integral read without integrating to the
endpoint.

\subsection*{L4$'$ --- sticking and release for prescribed forcing}

\begin{lemma}[sticking and release for prescribed $R$]\label{lem:release}
Let $R\in C^1$ be prescribed and satisfy \textup{(H0$'$)--(H2$'$)}, and
$\Ar(t)=\int_{\thit}^tR$. Then
\begin{enumerate}
\item[(i)] $\Psi(u(t))=\Psi(u_0)-c\beta\Ar(t)+\mathrm{corr}(t)$ with
$0\le\mathrm{corr}(t)=\int c\log(1+w)ds=O(\log\beta/\beta)$; on compacts of
$(\thit,\texit)$, $u=c\beta|\Ar(t)|(1+O(\log\beta/(\beta a_0)))$, whence
$z\sim-\beta^{-1}\log\beta$;
\item[(ii)] $\trel-\texit=\dfrac{\mathrm{corr}(\beta)}{c\beta R(\texit)}(1+o(1))
=O(\log\beta/\beta^2)$;
\item[(iii)] a perturbation $\tilde R=R+\delta$ shifts $\texit$ by at most
$\|\delta\|_\infty(\texit-\thit)/R(\texit)+O(\|\delta\|_\infty^2)$.
\end{enumerate}
No division by $R$ occurs: the layer is held by the area, not by the sign of $R$.
\end{lemma}

\noindent\emph{Provenance: identical to Lemma~4 of the paper --- which is already
stated for prescribed $R$ and never mentions $B$.} Step~1 is the Chaplygin
comparison, valid because the right-hand side of
$\dot\Psi=-c\beta R+c\log(1+w(\Psi))$ is decreasing in $\Psi$; steps~2--4 use
(H1$'$)--(H2$'$) and the implicit function theorem at the transversal zero
$\Ar(\texit)=0$. In particular the lemma imposes \emph{no} sign structure on $R$
beyond the two transversal zeros of $\Ar$: $R$ may change sign any number of
times inside the episode. This is the point taken up in \cref{rem:numerics}.

\subsection*{L5$'$ --- coupled bootstrap by continuous induction}

\begin{lemma}[coupled bootstrap]\label{lem:bootstrap}
Under \textup{(H0$'$)--(H5$'$)} there is $\beta_0<\infty$ such that for
$\beta\ge\beta_0$ the coupled system satisfies the two-term law
\begin{equation}\label{eq:twoterm}
t^\beta_{\mathrm{release}}
=\texit-\frac{1}{cR^{\hyb}(\texit)}\cdot\frac1\beta
+\frac{\mathrm{corr}(\beta)}{c\beta R^{\hyb}(\texit)}+o(\beta^{-1}),
\qquad \mathrm{corr}(\beta)=O(\log\beta/\beta),
\end{equation}
with $M_2=C_fC_4/c$ independent of the continuation hypothesis.
\end{lemma}

\noindent\emph{Provenance: structure of Lemma~5 of the paper; the continuation is
rewritten in full because it is the step where the lemmas are mutually
dependent, and because step~3 must now consume \eqref{eq:Rpert} through
(H3$'$) and the new live-unit term instead of ``$R$ is affine in $\theta_B$''.}

\begin{proof}
\emph{Step 0 (the dependency cycle, and how it is broken).} The three lemmas are
mutually dependent: \cref{lem:gauge}'s bound on $\int c\log(1+w)$ presupposes the
depth law of \cref{lem:release}(i), which presupposes the forcing actually felt
by $A$, which is $R^\beta$ and is controlled by \cref{lem:gronwall}, which
presupposes the layer bound of \cref{lem:gauge}. The cycle is not vicious because
every implication is \emph{causal in time}: each estimate on $[\thit,s]$ uses only
data on $[\thit,s]$. We therefore run a continuous induction in $s$.

\emph{Step 1 (the induction hypothesis).} Fix
$\varepsilon_\beta=C_\varepsilon\log\beta/\beta$ with
$C_\varepsilon=4C_4/(cR^{\hyb}(\texit))$ and a collar constant
$C_{\mathrm{collar}}$. For
$s\in I:=[\tcross+C_{\mathrm{collar}}/\beta,\ \texit-\varepsilon_\beta]$ let
$\mathcal P_s$ be the conjunction of
\begin{align}
&\text{(a)}\quad \sup_{[0,s]}\|\theta_B^\beta-\theta_B^{\hyb}\|
\ \le\ 2M_2\frac{\log\beta}{\beta^2},
\tag{a}\label{indA}\\
&\text{(b)}\quad u(t)\ \ge\ u_0\quad\text{for all }
t\in[\tcross+C_{\mathrm{collar}}/\beta,\ s],
\tag{b}\label{indB}\\
&\text{(c)}\quad \int_{\thit}^{s}c\log(1+w)\,dt\ \le\ 2C_4\frac{\log\beta}{\beta}.
\tag{c}\label{indC}
\end{align}
The collar is necessary: at $\tcross$ the margin crosses $\beta^{-1}\log w_0$
continuously, so the layer can be strict only after an $O(\beta^{-1})$ window.
Let $S=\{s\in I:\mathcal P_s\ \text{holds}\}$.

\emph{Step 2 ($S\neq\emptyset$).} \Cref{lem:corner-in} supplies the entry datum
with slack: $\|\theta^\beta(\tcross)-\theta^\infty(\thit)\|=O(\beta^{-1})$, the
layer is entered transversally with speed $v$, and by \cref{lem:gronwall}
step~4 the $B$-gap at $\thit$ is $O(\beta^{-2})$, an order below \eqref{indA}.
On a collar of width $C_{\mathrm{collar}}/\beta$ the integrand of \eqref{indC} is
$O(1)$, so \eqref{indC} holds with room. Hence $\mathcal P_s$ holds for $s$ just
past the collar and $S$ contains a right-neighbourhood of the left endpoint of
$I$.

\emph{Step 3 ($S$ closed in $I$).} All three quantities in
\eqref{indA}--\eqref{indC} are continuous, indeed monotone non-decreasing, in
$s$: $\sup$ and $\int$ over $[0,s]$, and \eqref{indB} is a closed condition on a
closed interval. Non-strict inequalities pass to limits, so $S$ is closed.

\emph{Step 4 ($S$ open in $I$: the improvement).} Let $s\in S$. We show the three
bounds hold on $[0,s]$ with the constants \emph{halved}; openness then follows
from continuity. Assume $\mathcal P_s$.
\begin{itemize}
\item[(4.1)] From \eqref{indC} and \cref{lem:gronwall} (whose hypothesis is
exactly \eqref{indC}) we get
$\sup_{[0,s]}\|\Delta\|\le M_2\log\beta/\beta^2+O(\beta^{-2})+
H(C_3e^{-\beta\delta_1}+C_5\sqrt mC_{\mathrm{res}}e^{-\beta\delta_B})$. For
$\beta\ge\beta_0$ with $\beta_0$ large enough that the last two groups are below
$\tfrac12M_2\log\beta/\beta^2$, this is
$\le\tfrac32M_2\log\beta/\beta^2<2M_2\log\beta/\beta^2$: strict improvement of
\eqref{indA}. \emph{The improvement is unconditional}: $M_2=C_fC_4/c$ is built
from $C_f$ (a property of $A$ and of the features), $C_4$ (a property of
$R^{\hyb}$), and $c$; the factor $2$ in the hypothesis never enters $M_2$. This
is what makes the induction non-circular.
\item[(4.2)] By \eqref{eq:Rpert} the forcing felt by $A$ on $[0,s]$ is
$R^\beta=R^{\hyb}+\delta$ with
$\|\delta\|_\infty\le2\mathrm{Lip}_RM_2\log\beta/\beta^2+m\beta^{-1}e^{-\beta\delta_B}$.
Enlarging $\beta_0$ so that $\|\delta\|_\infty H\le a_0/2$, the perturbed area
satisfies $|\Ar^\beta(t)|\ge\tfrac12|\Ar^{\hyb}(t)|$ on
$[\thit+\delta,\texit-\varepsilon_\beta]$ and the two transversal zeros move by
$O(\log\beta/\beta^2)$ by \cref{lem:release}(iii). Feeding this perturbed forcing
into \cref{lem:release}(i) gives, on
$[\tcross+C_{\mathrm{collar}}/\beta,\ s]$,
\[
u(t)\ \ge\ \tfrac12c\beta\bigl(|\Ar^{\hyb}(t)|-\|\delta\|_\infty H\bigr)\ >\ u_0 ,
\]
the last inequality because $|\Ar^{\hyb}(t)|\ge\kappa\,C_{\mathrm{collar}}/\beta$
just past the collar (transversality at $\thit$) and
$|\Ar^{\hyb}(t)|\ge\kappa\varepsilon_\beta=\Theta(\log\beta/\beta)$ just before
$\texit-\varepsilon_\beta$ (transversality at $\texit$), while
$\|\delta\|_\infty H=O(\log\beta/\beta^2)$. Choosing $C_{\mathrm{collar}}$ and
$C_\varepsilon$ as above makes the inequality strict with a fixed margin: strict
improvement of \eqref{indB}.
\item[(4.3)] With \eqref{indB} strict on $[0,s]$, the three-window computation of
\cref{lem:gauge} step~4 applies with $R^\beta$ in place of $R$; by (4.2) its two
transversality constants and its bulk constant $a_0$ are within a factor $2$ of
the hybrid ones, so it returns
$\int_{\thit}^{s}c\log(1+w)\le\tfrac32 C_4\log\beta/\beta$ (the factor $\tfrac32$
absorbing the factor-$2$ degradation of the constants against the slack we left):
strict improvement of \eqref{indC}.
\end{itemize}
Since \eqref{indA}--\eqref{indC} hold on $[0,s]$ with strict inequality and all
three functionals are continuous in $s$, they persist on $[0,s+\eta)$ for some
$\eta>0$; hence $S$ is open in $I$.

\emph{Step 5 (conclusion of the induction).} $I$ is a non-degenerate interval,
hence connected; $S\subseteq I$ is non-empty, open and closed; therefore $S=I$.
Equivalently: $s^\ast=\sup S$ cannot be interior, since at an interior $s^\ast$
continuity plus the strict improvement of step~4 would extend $S$ past $s^\ast$.
All three bounds therefore hold on the whole episode up to
$\texit-\varepsilon_\beta$.

\emph{Step 6 (terminal window and the two-term law).} On
$[\texit-\varepsilon_\beta,\trel]$ apply \cref{lem:release}(ii)--(iii) with the
forcing frozen at $R^{\hyb}+O(\log\beta/\beta^2)$. Three contributions add:
\begin{enumerate}
\item the displacement caused by $\delta$, $O(\log\beta/\beta^2)$ by (iii);
\item the correction delay $+\mathrm{corr}(\beta)/(c\beta R^{\hyb}(\texit))$ by
(ii);
\item the area cut at entry. By \cref{lem:corner-in} the episode starts with a
delay $K/\beta$, $K=1/(c|R^{\hyb}(\thit)|)$, during which the forcing is
$R^{\hyb}(\thit)<0$, so an area $\approx R^{\hyb}(\thit)K/\beta$ is removed from
$\Ar$; by transversality at $\texit$ this advances the exit by
\[
-\frac{K\,|R^{\hyb}(\thit)|}{R^{\hyb}(\texit)\,\beta}
=-\frac{1}{c\,R^{\hyb}(\texit)\,\beta},
\]
the arrival speed cancelling against $K$.
\end{enumerate}
Summing gives \eqref{eq:twoterm}.
\end{proof}

\begin{remark}[the cancellation is arithmetic in $R$]\label{rem:cancel}
Item~3 of step~6 is the entire content of the claim that the first-order
coefficient does not know about $B$. Both factors, $K=1/(c|R^{\hyb}(\thit)|)$
from \cref{lem:corner-in} and $|R^{\hyb}(\thit)|$ from the cut area, are written
in $R^{\hyb}$ and $c$ alone; their product is $1/c$, and the division by the exit
speed $cR^{\hyb}(\texit)$ leaves $-1/(cR^{\hyb}(\texit))$. Nothing about the live
subsystem --- not $m$, not $G_B$, not the readout scale --- survives to that
coefficient. Any $m$-dependence must therefore appear in the
\emph{remainder} $o(\beta^{-1})$, i.e. in $\beta_0$, and not in the law.
\end{remark}

\subsection*{L6$'$ --- exit corner, branch selection, scalar reduction}

\begin{lemma}[exit corner, branch selection and scalar reduction]\label{lem:corner-out}
Under \textup{(H0$'$)--(H5$'$)}, after release: (i) the $\beta$-trajectory
crosses the exit corner in time $\tfrac{1}{\beta v'}(1+o(1))$ with
$v'=cR^{\hyb}(\texit)>0$; (ii) the ReLU flow restarted at
$(\theta_A(\thit),\theta_B^{\hyb}(\texit))$ has forward uniqueness on the
activating branch, and that branch is selected by the limit itself, the frozen
Filippov branch being incompatible with release at $\texit+O(\beta^{-1})$;
(iii) classical Gr\"onwall glues on $[\trel,H]$ at rate $O(\log\beta/\beta)$.
Under (H4$'$) the other data of $A$ contribute $O(e^{-\beta\delta_1})$, the
face-$j$ field is rank one along $(x_j,1)$, and the episode reduces to the scalar
problem.
\end{lemma}

\noindent\emph{Provenance: identical to Lemma~6 of the paper with $R$ generic,
under the reinforced (H4$'$).} (i) mirrors \cref{lem:corner-in} with $v'$;
(ii) needs the active field to be locally Lipschitz on the side $z\ge0$ at the
restart point, which with $m$ live units holds provided no live unit sits exactly
on its kink at $\texit$ --- this is precisely the uniform margin $\delta_B$ of
(H4$'$), and it is the second place (after step~2b of \cref{lem:gronwall}) where
the reinforcement is used; the Filippov alternative is excluded by
\cref{lem:bootstrap}, independently of $m$; (iii) is classical Gr\"onwall on the
final interval, where all margins are bounded away from zero. Step~4, the scalar
reduction, uses \emph{only} that the other margins \emph{of $A$} are
$\le-\delta_1$ --- a statement about $A$, not about $B$ --- and is therefore
verbatim for every $m$.

\section{Proof of the theorem}\label{sec:proof}

\begin{proof}[Proof of \cref{thm:forcing}]
By \cref{lem:corner-out} step~4 the field on $\theta_A$ splits, under (H4$'$),
into a rank-one part along $(x_j,1)$ and terms of size $O(e^{-\beta\delta_1})$:
the tangential components of $\theta_A$ are frozen at an exponential rate and
the normal component is exactly the scalar problem of
\crefrange{lem:gauge}{lem:bootstrap}. This is what licenses treating the episode
as one-dimensional in $z$ for any $m$.

On $[0,\thit-\delta]$ both flows are affine with the same PSD Gram, and step~4 of
the proof of \cref{lem:gronwall} bounds their gap by $O(\beta^{-1})$ in the
$A$-block and $O(\beta^{-2})$ in the $B$-block. \Cref{lem:corner-in} then
delivers the entry corner: the crossing time is $\thit+1/(\beta v)(1+o(1))$ and
the state at crossing is $O(\beta^{-1})$ from $\theta^\infty(\thit)$.

On $[\thit+\delta,\texit-\delta]$, \cref{lem:bootstrap} closes the loop: its
continuous induction (steps~1--5) maintains simultaneously the layer, the bound
$\|\theta_B^\beta-\theta_B^{\hyb}\|\le2M_2\log\beta/\beta^2$ and the correction
integral. Inside the layer \cref{lem:release}(i) gives
$u=c\beta|\Ar^{\hyb}|(1+o(1))$, i.e. the sticking depth
$z=-\beta^{-1}\log\beta(1+o(1))$: this is item~(i) of the theorem, and it is the
largest of the three contributions to \eqref{eq:sup}, the other two being
$O(\log\beta/\beta^2)$ for $\theta_B$ (item~(iv)) and $O(e^{-\beta\delta_1})$ for
the tangential part of $\theta_A$. Hence the supremum over this stage is
$O_\delta(\log\beta/\beta)$, carried by the normal coordinate alone.

The terminal window and the release law are step~6 of \cref{lem:bootstrap}: the
correction delay of \cref{lem:release}(ii) is $O(\log\beta/\beta^2)$, the
perturbation of the forcing shifts $\texit$ by $O(\log\beta/\beta^2)$ by
\cref{lem:release}(iii), and the area cut by the entry corner advances the exit
by $-1/(cR^{\hyb}(\texit)\beta)$. Summing gives \eqref{eq:twoterm}, hence
item~(ii), \eqref{eq:release}.

Finally \cref{lem:corner-out} crosses the exit corner in time $O(\beta^{-1})$,
selects the activating branch, and glues the two flows on $[\trel,H]$ at rate
$O(\log\beta/\beta)$. Taking the maximum of the four stage bounds gives
\eqref{eq:sup}, item~(iii); item~(iv) is \cref{lem:gronwall} combined with the
induction. The threshold $\beta_0$ is the largest of the thresholds required in
steps (4.1)--(4.3) of \cref{lem:bootstrap} and of $\beta_1$ of
\cref{lem:gronwall} step~2b.
\end{proof}

\section{Uniformity in the size of the live subsystem}\label{sec:uniform}

\begin{remark}[what depends on $m$, and how]\label{rem:uniform}
We collect the $m$-dependence of every constant that enters
\cref{thm:forcing}. Write $X_{\max}=\max_i\|(x_i,1)\|$ and
$\|a\|_2=(\sum_{k\ge1}a_k^2)^{1/2}$.
\begin{enumerate}
\item[(U1)] \textbf{The first-order coefficient is exactly $m$-free.} By
\cref{rem:cancel}, $-1/(cR^{\hyb}(\texit))$ is built from $c$ (a property of the
datum $j$ and of $n$) and from the hybrid forcing at one instant. The cancellation
that produces it is arithmetic in $R$. Structurally, $m$ cannot appear.
\item[(U2)] $C_4=2/(c\min(|R^{\hyb}(\thit)|,R^{\hyb}(\texit)))+\Tep/a_0$ is
written purely in $R^{\hyb}$ and $c$: $m$-free \emph{given the forcing}. The
forcing of course depends on $m$; the point is that $C_4$ is a functional of the
forcing only, so all $m$-dependence is routed through the observable quantities
$R^{\hyb}(\thit)$, $R^{\hyb}(\texit)$, $a_0$, $\Tep$.
\item[(U3)] $C_f=\|\varphi_j\|/n\le X_{\max}\|a\|_2/n$. With unit readouts
$a_k\equiv1$ this is $\le X_{\max}\sqrt m/n$; with NTK readouts
$a_k=m^{-1/2}$ it is $O(1)$. \emph{At worst $C_f\lesssim\sqrt m$.}
\item[(U4)] $M_2=C_fC_4/c\lesssim\sqrt m$. This is the substantive gain of
\cref{lem:gronwall}: the published constant is $e^{L_BH}C_fC_4/c$ with
$L_B=\lambda_{\max}(G_B)$, and $\lambda_{\max}(G_B)\sim m$ for unit readouts, so
the published $M_2$ grows like $\sqrt m\,e^{\Theta(mH)}$. \textbf{There is no
$m\cdot H$ restriction} in the PSD version: the horizon enters only through
$\Tep/a_0$ inside $C_4$ and through $H$ multiplying quantities that are beyond
all orders in $\beta$.
\item[(U5)] The live-unit term \eqref{eq:GS} carries $\sqrt m\,C_{\mathrm{res}}$
(crudely $mn$ scalar discrepancies) times $e^{-\beta\delta_B}$. It affects only
$\beta_1=\delta_B^{-1}\log(m^2H\beta^2)$, i.e. $\beta_0$ grows
\emph{logarithmically} in $m$ through this channel.
\item[(U6)] What is \emph{not} controlled: $\beta_0$ itself. The remaining
$m$-dependence of the theorem lives entirely in $\beta_0$, and we do not exhibit
it --- the same gap as in the published two-neuron theorem. See \cref{sec:open}.
\end{enumerate}
Conclusion: under (H0$'$)--(H5$'$) every explicit constant of
\cref{thm:forcing} depends on $m$ at most polynomially (indeed at most like
$\sqrt m$), and the law itself not at all.
\end{remark}

\section{Scope of the numerical verification}\label{sec:numerics-scope}

\begin{remark}[what the numerics cover and what they do not]\label{rem:numerics}
Three independent numerical campaigns bear on this theorem; none of them is part
of its proof.
\begin{enumerate}
\item \textbf{The $m$-sweep with NTK readout} ($66$ instances passing
(H1$'$)--(H4$'$): $23$ at $m=4$, $21$ at $m=8$, $22$ at $m=16$, plus a
$57$-instance two-neuron baseline; $\beta=10^4$, explicit Euler with
$dt_0=10^{-5}$ on the hybrid and $dt=2\cdot10^{-6}$ on the coupled flow,
$H=1.8$). No free parameter is fitted: the predicted coefficients spanned
$[-109,-14]$, a factor of $8$, and the measured ones $[-132,-14]$, with
$R^2=0.950$ and a regression slope of $1.154$. Median relative error
$0.80\%/6.85\%/5.58\%$ at $m=4/8/16$ against $2.68\%$ for the two-neuron
baseline. Restricting to a matched band $R^{\hyb}(\texit)\in[0.05,0.13]$ gives
$1.81\%$ at $m=4$ against $1.79\%$ at the baseline, but still $7.40\%$ at $m=8$;
a discretization control at $dt_0/5$ leaves the residual unchanged, so it is not
the integrator. In an OLS of the error on $(\log_2m,\ 1/R^{\hyb}(\texit),\
\Tep)$, the coefficient of $\log_2m$ is not significant ($t=0.71$) while
$1/R^{\hyb}(\texit)$ is ($t=2.56$) --- consistent with the theorem's remainder
being governed by exit transversality rather than by $m$, but not conclusive at
$n=66$.
\item \textbf{Bit-identical cloning} ($23$ instances). Each accepted $m=4$
instance is replayed with every live unit split into $p$ copies carrying readout
$a_k/\sqrt p$ and parameters $(w_k,b_k)/\sqrt p$, so that $f_B$, the NTK and the
whole flow --- hence $R(t)$ --- are unchanged while $m\to pm$. The measured
errors are identical to machine precision at $m=4,8,16$ (median $0.80\%$, max
$16.17\%$ in all three columns), and $\lambda_{\max}(G_B)$ is likewise invariant.
This isolates the claim of (U1): \emph{at fixed forcing, the number of live units
does nothing}. It is the closest numerical analogue of \cref{rem:cancel}.
\item \textbf{The unit-readout family} ($67$ instances: $22/24/21$ at
$m=4/8/16$, median error $1.16\%/4.52\%/5.35\%$, max $\approx28\%$). This is the
family of \eqref{eq:net} with $a_k\equiv1$, in which $\lambda_{\max}(G_B)$ grows
with $m$ and the published $M_2$ would be exponential. The law survives, with the
same qualitative degradation pattern; under the PSD constant \eqref{eq:M2} this
is expected, under $e^{L_BH}C_fC_4/c$ it would not be. Together with the previous
item this is the empirical case for \cref{lem:gronwall} being the right estimate,
and it is why the exponential was worth removing.
\item \textbf{The $\beta$ sweep on the worst instances.} The six worst $m=8$
instances of item~1 (errors $11.9$--$20.5\%$ at $\beta=10^4$) were rerun at
$\beta=3\cdot10^4$ with the same integrator and the same predicted coefficients,
which are properties of the hybrid reference and therefore unchanged. The errors
collapse to
\[
20.45\to0.35,\quad 20.20\to4.93,\quad 17.30\to0.20,\quad
13.52\to4.23,\quad 12.62\to0.30,\quad 11.88\to0.50\ \ (\%),
\]
i.e. by roughly the factor $3$ in $\beta$ expected of an $O(\beta^{-1})$
remainder relative to the $O(\beta^{-1})$ leading term, on every instance.
This is the sharpest available evidence that the $m=8,16$ degradation of item~1
is a $\beta_0$ effect --- the asymptotic regime bites later --- and not a failure
of the law; see (O4).
\end{enumerate}
\textbf{The gap between the numerics and the theorem.} In every instance
generated by these designs, $R^{\hyb}$ changes sign \emph{once} inside the
episode, because the acceptance filter selects a single entry--exit episode with
$\Ar^{\hyb}<0$ throughout. \Cref{lem:release} imposes no such restriction: the
Chaplygin comparison of its step~1 needs only monotonicity of the right-hand side
in $\Psi$, and (H2$'$) constrains the zeros of the \emph{area}, not of $R$. The
theorem therefore covers forcings that oscillate arbitrarily often inside the
episode, and the numerics do not test that regime. Conversely, the acceptance
rate of (H1$'$)--(H2$'$) is about $3\%$ of random initializations at every $m$,
so the hypotheses are verifiable a posteriori by one hybrid integration but are
not guaranteed a priori by any initialization criterion we know.
\end{remark}

\section{Steps we did not close}\label{sec:open}

Stated plainly, and separated from what is proved above.

\begin{enumerate}
\item[(O1)] \textbf{$\beta_0$ is not exhibited.} Steps (4.1)--(4.3) of
\cref{lem:bootstrap} require $\beta$ large enough for three explicit
inequalities, and $\beta_1$ of \cref{lem:gronwall} step~2b for a fourth. Each is
traceable through the chain of constants, but we have not composed them into a
number. This gap is inherited verbatim from the published two-neuron theorem; it
is the only place where an $m$-dependence beyond $\sqrt m$ could still hide, and
therefore the honest location of the residual risk flagged in \cref{rem:uniform}
(U6).
\item[(O2)] \textbf{The piecewise version of (H4$'$) is not done.} With $m$ live
units, $R^{\hyb}$ is analytic between switching events and has a corner at each
event, so it is $C^0$ and piecewise $C^1$, not $C^1$. (H4$'$) removes all corners
by decree. The proofs above use $C^1$ of $R$ only through
$\|\dot R\|_\infty\le C_R$ (which survives corners; Lipschitz suffices) and
through transversality of the zeros of $\Ar$ at $\thit$ and $\texit$. A weaker
hypothesis (H4$''$) --- \emph{no live unit switches in a fixed neighbourhood of
$\thit$ and of $\texit$}, with $R$ merely Lipschitz in between --- ought to
suffice, but we have not done that bookkeeping, and (H4$'$) is additionally used
for branch selection in \cref{lem:corner-out}(ii) and for the uniform live
margins in \cref{lem:gronwall} step~2b, where the weakening is \emph{not}
obviously free.
\item[(O3)] \textbf{The $o(1)$ of \cref{lem:corner-in} is not quantified.} The
entry-corner estimate carries $O(\delta_z)+o(1)$ and we do not sharpen it to
$O(\log\beta/\beta)$; consequently the remainder of \eqref{eq:release} is
$o(\beta^{-1})$ and not the $O(\log\beta/\beta^2)$ that the numerics suggest.
Inherited from the paper.
\item[(O4)] \textbf{The measured degradation at $m=8,16$ is not explained by the
constants.} \Cref{rem:uniform} predicts $M_2\lesssim\sqrt m$, i.e. a mild growth,
whereas the median error grows by a factor $4$--$8$ from $m=4$ to $m=8$ in the
NTK family. Three pieces of evidence say the discrepancy lives in $\beta_0$ and
not in the law --- the cloning experiment (item~2 of \cref{rem:numerics}) shows
that at fixed forcing $m$ does nothing; the regression shows the error tracking
$1/R^{\hyb}(\texit)$; and the $\beta$ sweep (item~4) drives the six worst
instances from $12$--$20\%$ to $0.2$--$4.9\%$ by tripling $\beta$. But a bound
that \emph{predicts} the observed degradation from exit transversality alone is
not in this note, and without (O1) we cannot certify that $\beta_0$ stays
polynomial in $m$. Until then, the uniformity claim of \cref{rem:uniform} is a
statement about the constants we wrote, not a demonstrated property of the error.
\item[(O5)] \textbf{Step~4 of \cref{lem:gronwall} is new and unaudited.} The
pre-hit initial gap $\rho(\thit)=O(\beta^{-2})$ rests on splitting the drift into
$-G_B\Delta-G_{BA}\delta\theta_A$ and on the pointwise weighted bound
\eqref{eq:prehitA}, which is step~1 of \cref{lem:corner-in} read before
integrating to the endpoint. The published proof leaves this initial gap
implicit. We believe the estimate; it has not been checked by a referee or by a
script.
\end{enumerate}

\bigskip
\noindent\textbf{Status.} Sections \ref{sec:setting}--\ref{sec:proof} are a
proof \emph{conditional} on (O1) and (O5) and \emph{restricted} by (O2).
\Cref{rem:numerics} is a report of executed runs, not evidence for the proof.
\Cref{rem:uniform} is proved for the constants displayed and open for $\beta_0$.

\end{document}
